\documentclass[11pt,a4paper]{article}
% Mise en page commune aux guides MyAnimLab (utilisateur, administrateur).
\usepackage[utf8]{inputenc}
\usepackage[T1]{fontenc}
\usepackage[french]{babel}
\usepackage[margin=2.2cm]{geometry}
\usepackage{enumitem,tabularx,booktabs}
\usepackage[table]{xcolor}
\usepackage{xurl}
\usepackage[hidelinks]{hyperref}
\usepackage{fancyhdr}
\usepackage[most]{tcolorbox}
\sloppy

\definecolor{accent}{HTML}{1F5F8B}
\definecolor{ok}{HTML}{2E7D32}
\definecolor{warn}{HTML}{C76A00}

% Chemin de menu : \menu{Animaux > Tournée quotidienne}
\newcommand{\menu}[1]{\textsf{\textbf{\color{accent}#1}}}
% Bouton ou libellé de l'interface
\newcommand{\btn}[1]{\textsf{«~#1~»}}
% Page ou fichier technique
\newcommand{\f}[1]{\nolinkurl{#1}}

% Procédure pas à pas
\newtcolorbox{procedure}[1]{colback=accent!4,colframe=accent,fonttitle=\bfseries\sffamily,
  title={#1},left=2mm,right=2mm,top=1mm,bottom=1mm,breakable}
% Règle à connaître
\newtcolorbox{regle}{colback=warn!6,colframe=warn,fonttitle=\bfseries\sffamily,
  title={À savoir},left=2mm,right=2mm,top=1mm,bottom=1mm,breakable}
% Traçabilité réglementaire
\newtcolorbox{trace}{colback=ok!6,colframe=ok,fonttitle=\bfseries\sffamily,
  title={Traçabilité},left=2mm,right=2mm,top=1mm,bottom=1mm,breakable}

\setlist[enumerate]{nosep,leftmargin=*}
\setlist[itemize]{nosep,leftmargin=*}

\pagestyle{fancy}
\fancyhf{}
\cfoot{\thepage}

\lhead{MyAnimLab --- Guide de l'utilisateur}
\rhead{Version du 18 octobre 2026}

\title{\textbf{MyAnimLab}\\[4pt]\large Guide de l'utilisateur}
\author{Animaliers, expérimentateurs, techniciens}
\date{18 octobre 2026}

\begin{document}
\maketitle
\thispagestyle{fancy}

\noindent Ce guide décrit les tâches quotidiennes dans MyAnimLab : tenir les
cages, faire la tournée, suivre les animaux, la reproduction et les
expérimentations, administrer les traitements. Les actions réservées aux
administrateurs (comptes, paramètres, autorisations de projet, corrections du
registre) sont décrites dans le \emph{Guide de l'administrateur}.

\tableofcontents
\newpage

\section{Premiers pas}

\subsection{Se connecter}
Les comptes sont créés par un administrateur : il n'y a pas d'inscription
libre. Votre identifiant est votre adresse e-mail.

\begin{procedure}{Première connexion}
\begin{enumerate}
  \item Ouvrez l'adresse de l'application et saisissez l'e-mail et le mot de
    passe initial transmis par l'administrateur.
  \item Si l'établissement l'exige pour votre rôle, configurez la double
    authentification (section~\ref{sec:2fa}) : aucune autre page n'est
    accessible avant.
  \item Changez votre mot de passe dans \menu{Paramètres > Profil utilisateur}.
\end{enumerate}
\end{procedure}
\capture[0.62\linewidth]{u_connexion}{Page de connexion.}

\begin{regle}
\begin{itemize}
  \item Après 3 mots de passe erronés, le compte est verrouillé 15 minutes.
  \item Mot de passe oublié : lien \btn{Mot de passe oublié} de la page de
    connexion. Un e-mail de réinitialisation est envoyé.
  \item Sans activité pendant la durée fixée par l'administrateur (30 minutes
    par défaut), la session se ferme. Il faut alors se reconnecter.
\end{itemize}
\end{regle}

\subsection{Double authentification}\label{sec:2fa}
\begin{procedure}{Activer la double authentification}
\begin{enumerate}
  \item \menu{Paramètres > Double authentification}.
  \item Scannez le QR code avec une application d'authentification
    (Google Authenticator, Microsoft Authenticator, FreeOTP…).
  \item Saisissez le code à 6 chiffres affiché par l'application.
  \item \textbf{Conservez les dix codes de secours} affichés : chacun sert
    une seule fois si vous perdez votre téléphone.
\end{enumerate}
\end{procedure}
À chaque connexion, le code à 6 chiffres est demandé après le mot de passe.
Après 5 codes erronés, le compte est verrouillé 15 minutes. En cas de perte du
téléphone et des codes de secours, demandez une réinitialisation à un
administrateur.

\subsection{Se repérer}
\begin{itemize}
  \item \menu{Accueil} : statistiques des cages, alertes du jour (actions
    de point limite, portées à sevrer, en rouge si le sevrage est en retard,
    avancement de la tournée) et \textbf{notes} (pense-bêtes partagés).
  \item \menu{Tableaux de bord} : cages d'hébergement, d'accouplement,
    d'expérimentation, toutes les cages par salle.
  \item \menu{Animaux} : tournée quotidienne, scan QR, recherche, transferts,
    créations en lot, essais et expériences, grilles cliniques, sanitaire,
    terrariums, registre des entrées et sorties, pharmacie.
  \item La \textbf{cloche} affiche vos notifications : rappels d'actes,
    portées à sevrer, alertes de point limite, validations de formation…
  \item Le champ de recherche de l'en-tête trouve un animal par son numéro ou
    sa cage.
  \item La langue (français ou anglais) se choisit dans l'en-tête ou dans le
    profil.
\end{itemize}
\capture[0.85\linewidth]{u_accueil}{Accueil : action de point limite requise, portée dont le sevrage est en retard, avancement de la tournée.}

\section{Les cages}

\subsection{Tableaux de bord et fiche de cage}
Chaque catégorie de cages a son tableau de bord : salle, rack, animaux
vivants, recherche, filtre par salle. Une sélection groupée permet de changer
de salle ou de rack, ou de passer une cage entre hébergement et
expérimentation. Un clic sur une cage ouvre sa \textbf{fiche}, qui présente :
\begin{itemize}
  \item l'emplacement (salle, rack) et le statut ;
  \item les animaux vivants, avec leurs notes de santé ;
  \item pour une cage d'accouplement, l'encadré \btn{Portées en cours}
    (section~\ref{sec:portees}) ;
  \item les \textbf{fichiers} joints et le journal de maintenance.
\end{itemize}

\subsection{Créer des cages et des animaux}
\begin{procedure}{Cage d'hébergement avec ses animaux}
\begin{enumerate}
  \item \menu{Tableaux de bord > Cages d'hébergement}, puis \btn{Ajouter une
    cage d'hébergement}.
  \item Choisissez la lignée, le nombre de mâles et de femelles et la date
    de naissance.
  \item Renseignez l'\textbf{entrée dans l'établissement} : provenance
    (fournisseur ou éleveur agréé, autre établissement utilisateur, autre),
    fournisseur, référence du bon de livraison et date d'entrée. Ces
    informations sont \textbf{obligatoires} pour le registre.
  \item Validez. Le numéro de cage et les numéros des animaux (par exemple
    \texttt{RSD-0002}) sont attribués automatiquement.
\end{enumerate}
\end{procedure}
Pour un arrivage important, \menu{Animaux > Créer des animaux et cages en
lot} répartit automatiquement les animaux sur le nombre minimal de cages,
selon la capacité par cage indiquée.

\begin{procedure}{Cage d'accouplement}
\begin{enumerate}
  \item \menu{Tableaux de bord > Cages d'accouplement}, puis \btn{Ajouter une
    cage d'accouplement}.
  \item Indiquez le croisement, le responsable, l'autorisation de projet, la
    salle et le rack, ainsi que les utilisateurs de la cage.
  \item Choisissez les \textbf{reproducteurs existants} : au moins un mâle et
    une femelle. Ils sont transférés dans la nouvelle cage, et le mouvement est
    conservé dans leur historique.
  \item Les naissances s'enregistrent ensuite depuis la fiche de la cage
    (section~\ref{sec:portees}).
\end{enumerate}
\end{procedure}

\subsection{Entretien, notes et capteurs}
\begin{itemize}
  \item \textbf{Entretien} : depuis la fiche d'une cage (icône clé du journal
    de maintenance), \btn{Ajouter une entrée de maintenance}. Choisissez une ou
    plusieurs cages et ajoutez un commentaire facultatif (changement de
    litière, de biberon\ldots). Le journal de maintenance de chaque cage
    conserve la date et l'auteur.
  \item \textbf{Notes} : pense-bêtes affichés sur l'accueil (notes générales)
    ou sur la fiche d'une cage (notes de la cage). Une note ne peut être
    modifiée ou supprimée que par son auteur ou par un administrateur.
  \item \textbf{Capteurs} : \menu{Tableaux de bord > Capteurs IoT} affiche la
    température, l'humidité, l'éclairage et la pression des salles, quand
    l'administrateur a configuré les capteurs.
\end{itemize}

\subsection{Étiquettes et QR codes}\label{sec:etiquettes}
\menu{Tableaux de bord} : sélectionnez des cages puis imprimez leurs
étiquettes (A6, quatre par feuille A4, même mise en page pour toutes les
catégories). Le QR code de l'étiquette ouvre le contrôle de la cage
(section~\ref{sec:tournee}).

L'étiquette d'une cage d'accouplement porte les reproducteurs et un tableau
\textbf{Portées --- à compléter au crayon}. Les portées en allaitement y sont
imprimées : date de naissance, nés vivants, morts déjà saisies, sevrage
prévu. Les lignes vides servent aux événements qui surviennent après
l'impression.
\begin{procedure}{Garder la fiche de cage à jour sans la réimprimer}
\begin{enumerate}
  \item Notez au crayon sur la fiche chaque naissance, mort (date et nombre)
    et sevrage. Le sevrage prévu est égal à la date de naissance plus l'âge
    au sevrage imprimé sur la fiche.
  \item Scannez le QR code : la page de la cage propose \btn{Enregistrer une
    naissance} et, pour chaque portée, \btn{Décès / sevrage}. Enregistrez ce
    que vous venez de noter.
  \item Cochez la case \btn{Saisi} de la ligne. Une ligne non cochée signale
    une donnée encore absente du logiciel.
\end{enumerate}
\end{procedure}
\captures{u_etiquette}{0.4\linewidth}{Étiquette d'une cage d'accouplement : tableau des portées à compléter au crayon.}{u_qr_accouplement}{0.27\linewidth}{Page ouverte par le QR code : report des portées.}

\subsection{Fichiers joints à une cage}
\begin{procedure}{Ajouter un fichier}
\begin{enumerate}
  \item Ouvrez la fiche de la cage, section \btn{Gérer les fichiers}.
  \item \btn{Ajouter un fichier} : PDF, document Office, texte ou image,
    10~Mo au plus. Validez avec \btn{Envoyer}.
\end{enumerate}
\end{procedure}
\capture[0.8\linewidth]{u_fichiers}{Section « Gérer les fichiers » de la fiche d’une cage.}
Seuls les utilisateurs affectés à la cage et les administrateurs ajoutent ou
suppriment des fichiers. Tout utilisateur connecté peut les télécharger. Un
fichier dont le contenu ne correspond pas à son extension est refusé. Chaque
ajout et chaque suppression sont enregistrés dans le journal d'activité.

\subsection{Retirer une cage}
Une cage ne peut être retirée qu'une fois que chacun de ses animaux vivants a
reçu un devenir : transfert dans une autre cage, ou sortie avec son code. La
page \btn{Retirer la cage} demande ce devenir animal par animal. L'opération
est réalisée en entier ou pas du tout. Une cage qui a encore des portées en
allaitement ne peut pas être retirée. Une cage vidée de son dernier animal
est retirée automatiquement. Une cage retirée n'est pas remise en service
(seul un administrateur peut le faire, dans un cas d'anomalie).

\section{La tournée quotidienne}\label{sec:tournee}
Chaque cage en service doit être vue au moins une fois par jour.

\begin{procedure}{Faire la tournée}
\begin{enumerate}
  \item \menu{Animaux > Tournée quotidienne} : les cages sont classées par
    salle.
  \item Pour chaque cage : \btn{RAS} (une à une ou en groupe), ou
    signalement d'un problème sanitaire sur l'un de ses animaux.
  \item Sur téléphone : \menu{Animaux > Scanner le QR de cage}, scannez
    l'étiquette, puis validez la cage en un geste ou signalez un problème.
\end{enumerate}
\end{procedure}
\capture[0.9\linewidth]{u_tournee}{Tournée quotidienne : avancement par salle, cages vues, bouton RAS.}

\begin{regle}
\begin{itemize}
  \item Une cage déjà vue dans la journée (suivi expérimental, signalement
    sanitaire, scan QR) apparaît comme vue.
  \item Une \textbf{cage en expérimentation} ne peut pas être validée
    \btn{RAS}. Elle n'est validée qu'une fois saisi le suivi du jour de
    chacun de ses animaux en expérience, et confirmés ses actes du jour.
  \item En fin de journée, à l'heure fixée par l'administrateur (18:00 par
    défaut), les administrateurs reçoivent la liste des cages non vues.
  \item Les jours passés se consultent en lecture seule.
\end{itemize}
\end{regle}

\section{Les animaux}

\subsection{Trouver un animal}
Utilisez \menu{Animaux > Rechercher un animal} ou la recherche de l'en-tête :
par numéro, cage, lignée, sexe ou statut. Un numéro exact ouvre directement la
fiche.

\subsection{La fiche animal}
La fiche présente l'identité de l'animal (lignée, sexe, date de naissance,
\textbf{signe distinctif} : boucle, tatouage, puce, pelage…), sa cage, sa
filiation, ses expériences et sa santé. Elle comporte aussi le cadre
\btn{Registre des entrées et sorties}.

\begin{procedure}{Compléter l'entrée d'un animal}
Si la provenance est marquée \btn{À compléter} : dépliez \btn{Compléter
l'entrée dans l'établissement}, renseignez la provenance, le fournisseur, la
référence et la date d'entrée, puis enregistrez.
\end{procedure}

\begin{trace}
Une fois l'entrée enregistrée, seul un administrateur peut la corriger, avec
un motif. Les valeurs remplacées restent visibles dans \btn{Corrections du
registre} sur la fiche et dans le dossier BEA.
\end{trace}
\capture[0.8\linewidth]{u_registre_animal}{Cadre « Registre des entrées et sorties » d’un animal, avec une correction de sortie.}

\subsection{Planifier les utilisations}
\menu{Animaux > Planification des utilisations} : pour un animal, notez
l'utilisation prévue (par exemple « prélèvement J28 »), la date prévue, puis la
date de réalisation, et cochez \btn{Déjà réalisée}. La page liste les utilisations
prévues et réalisées.

\subsection{Transférer un animal}
\menu{Animaux > Transférer des animaux}, ou bouton \btn{Changer de cage} sur
la fiche de la cage. Chaque mouvement est conservé avec sa date, son auteur
et son motif. Inclure un animal dans une expérience, ou le transférer dans
une cage en expérimentation, fait passer cette cage en expérimentation.

\subsection{Enregistrer une sortie}
\begin{procedure}{Sortie d'un animal}
\begin{enumerate}
  \item Depuis la fiche : \btn{Sortie (mise à mort / mort / export)}.
  \item Choisissez le \textbf{code de sortie} : mis à mort (gestion de
    population, en expérimentation), euthanasie, atteinte de point limite,
    retrouvé mort (en expérimentation ou non), export vers une autre
    animalerie (l'établissement de destination est alors obligatoire).
  \item Indiquez la date et les circonstances, puis confirmez.
\end{enumerate}
\end{procedure}
La fiche de l'animal reste consultable indéfiniment. Ses inclusions en cours
sont clôturées et ses actions de point limite en attente sont soldées. Une
sortie erronée ne peut être corrigée que par un administrateur, avec un
motif.

\subsection{Signaler un problème sanitaire}
\begin{procedure}{Signalement et suivi}
\begin{enumerate}
  \item \menu{Animaux > Signaler un problème sanitaire}, ou depuis la fiche
    de cage ou la tournée.
  \item Choisissez la catégorie (blessure, perte de poids, comportement,
    signes respiratoires ou digestifs, peau, tumeur, etc.), la gravité et la
    description. Les administrateurs sont prévenus.
  \item Les suivis s'ajoutent depuis le signalement : chaque suivi indique le
    nouveau statut (ouvert, en surveillance, résolu) et la gravité.
\end{enumerate}
\end{procedure}
\begin{trace}
Le \menu{Registre sanitaire} conserve tous les signalements et leurs suivis.
Aucune entrée ne peut être modifiée ni supprimée. Il s'exporte en CSV, une
ligne par suivi.
\end{trace}
\capture[0.85\linewidth]{u_sanitaire}{Un signalement sanitaire et son suivi.}

\section{La reproduction}\label{sec:portees}

\subsection{Mettre des animaux en accouplement}
\menu{Tableaux de bord > Cages d'accouplement}, puis \btn{Nouvelle cage
d'accouplement}. Les reproducteurs sont des animaux déjà présents dans
l'établissement : ils sont transférés dans la nouvelle cage avec tout leur
historique.

\begin{procedure}{Créer une cage d'accouplement}
\begin{enumerate}
  \item \btn{Responsable} : le chercheur responsable (PI) de la cage.
  \item \btn{Croisement} : par exemple « RSD × RSD » ou
    « C57BL/6 × Cre+ ». Il figure sur l'étiquette.
  \item \btn{Numéro éthique} : l'autorisation de projet dont relèvent les
    animaux, s'il y a lieu.
  \item \btn{Utilisateur} : les personnes qui s'occupent de la cage ; elles
    reçoivent les rappels de sevrage.
  \item \btn{Reproducteurs mâles existants} et \btn{Reproductrices
    existantes} : tapez le numéro ou la lignée, puis choisissez parmi les
    animaux actifs des autres cages.
  \item \btn{Salle} et \btn{Rack} : la salle sert à organiser la tournée
    quotidienne.
  \item Enregistrez. Le numéro de cage est attribué automatiquement.
    Imprimez l'étiquette : elle comporte le tableau des portées à compléter
    au crayon et le QR code (section~\ref{sec:etiquettes}).
\end{enumerate}
\end{procedure}
\capture[0.55\linewidth]{a_accouplement}{Formulaire « Nouvelle cage d’accouplement ».}

\subsection{Portées en cours}
En tête de la fiche d'une cage d'accouplement, l'encadré \btn{Portées en
cours} liste les portées non sevrées. Pour chacune, il indique la date de
naissance, l'âge en jours, les petits vivants (morts enregistrées déduites) et
la date de sevrage prévue, en rouge si elle est dépassée. Les petits non
sevrés ne sont pas encore identifiés et ne sont pas comptés parmi les animaux
vivants de la cage. Les portées sevrées ou clôturées sont regroupées plus bas
dans \btn{Portées précédentes}.

\capture[0.8\linewidth]{u_portees_en_cours}{Encadré « Portées en cours » : le sevrage en retard est en rouge.}

\subsection{Enregistrer une naissance}
\begin{procedure}{Déclarer une portée}
\begin{enumerate}
  \item Sur la fiche de la cage, ou après le scan de son QR code :
    \btn{Portées : naissance → sevrage}.
  \item Section \btn{Enregistrer une naissance} : \btn{Date de naissance},
    \btn{Date d’accouplement} si elle est connue, \btn{Mère} et \btn{Père}
    parmi les reproducteurs de la cage, \btn{Petits nés vivants},
    \btn{Petits nés morts} et \btn{Remarques}.
  \item \btn{Enregistrer la naissance}. La date de sevrage prévue est
    calculée automatiquement (âge réglé par l'administrateur, 21 jours par
    défaut), et la portée apparaît dans l'encadré \btn{Portées en cours}.
  \item Reportez aussi la portée au crayon sur l'étiquette de la cage.
\end{enumerate}
\end{procedure}
\capture[0.85\linewidth]{u_naissance}{Section « Enregistrer une naissance ».}

\subsection{Pendant l'allaitement}
Depuis la page de la portée (bouton \btn{Ouvrir} dans la liste des portées,
ou lien de l'encadré) :
\begin{itemize}
  \item \btn{Enregistrer des décès de petits} : date, nombre et détails
    (« Retrouvé mort, euthanasié… »), puis \btn{Enregistrer les décès}. Le
    nombre de \btn{Petits vivants actuellement} diminue d'autant.
  \item \btn{Sexage} : nombre de mâles et de femelles, dont la somme doit
    égaler le nombre de petits vivants, puis \btn{Enregistrer le sexage}.
    Ces nombres pré-remplissent le formulaire de sevrage.
  \item Le \btn{Journal} de la portée garde chaque événement avec sa date et
    son auteur.
\end{itemize}
\capture[0.85\linewidth]{u_portee}{Page d’une portée : résumé, actions et journal.}

\subsection{Sevrer une portée}
\begin{procedure}{Sevrer et identifier les petits}
\begin{enumerate}
  \item Depuis l'encadré \btn{Portées en cours} : \btn{Sevrer / mettre à
    jour}, puis la section \btn{Sevrer la portée}.
  \item \btn{Date de sevrage} et \btn{Lignée} des petits.
  \item Une ligne par groupe : \btn{Sexe}, \btn{Nombre} et
    \btn{Destination}, soit \btn{Nouvelle cage d’hébergement} (avec
    \btn{Salle} et \btn{Rack}), soit une cage existante. \btn{Ajouter un
    groupe (répartition dans plusieurs cages)} sert à répartir un même sexe
    dans plusieurs cages.
  \item Le compteur \btn{Sevrés : x / y vivants} vérifie que tous les petits
    sont comptés. S'il en manque, expliquez ce qui s'est passé dans
    \btn{Petits non sevrés : ce qui s’est passé}.
  \item \btn{Sevrer la portée}. Chaque petit devient un animal identifié :
    numéro de lignée, sexe, date de naissance, mère et père, et entrée au
    registre « né dans l'établissement ». Imprimez les étiquettes des
    nouvelles cages.
\end{enumerate}
\end{procedure}
\capture[0.85\linewidth]{u_sevrage}{Formulaire de sevrage : deux groupes vers de nouvelles cages d’hébergement.}

Une portée perdue se clôture sans sevrage : \btn{Clôturer la portée sans
sevrage}, avec la raison (« portée perdue, euthanasiée… »), puis
\btn{Clôturer la portée}.

\begin{regle}
Chaque matin, un rappel (notification et e-mail) signale aux utilisateurs de
la cage les portées à sevrer le jour même ou le lendemain. Un sevrage en
retard est aussi signalé aux administrateurs et aux responsables BEA. Une
portée sevrée ou clôturée ne peut plus être modifiée.
\end{regle}

\section{L'expérimentation}

\subsection{Organisation}
\menu{Animaux > Essais / protocoles / expériences}. Un \textbf{essai} est
rattaché à une autorisation de projet (numéro éthique) : sa période doit être
comprise dans la validité de l'autorisation. Un essai contient des
\textbf{plans / protocoles}, qui contiennent des \textbf{expériences}, qui
contiennent les \textbf{animaux inclus}.

\begin{procedure}{Créer un essai}
\begin{enumerate}
  \item \btn{Nouvel essai} : numéro, titre, \btn{Numéro éthique} (parmi les
    autorisations de projet enregistrées par l'administrateur), dates de
    début et de fin, statut et description.
  \item \btn{Créer l'essai}. La carte de l'essai rappelle l'autorisation,
    son numéro et sa validité (badge \btn{Valide} ou \btn{Expirée}).
  \item \btn{Modifier l’essai} corrige ces informations ; \btn{Rapport}
    imprime le rapport d'essai.
\end{enumerate}
\end{procedure}

\begin{procedure}{Préparer un plan / protocole}
\begin{enumerate}
  \item Dans l'essai, \btn{Nouveau plan / protocole} : numéro, titre, dates
    (comprises dans celles de l'essai), \btn{Sévérité attendue},
    \btn{Méthodologie} et \btn{Critères de point limite — description}
    (texte repris dans le rapport).
  \item Onglet \btn{Grille clinique} : choisissez la \textbf{grille
    d'évaluation clinique} du suivi quotidien, puis \btn{Enregistrer la
    grille}. Les grilles se créent dans \menu{Animaux > Grilles
    d'évaluation clinique} ; une grille déjà utilisée ne peut plus être
    modifiée : dupliquez-la pour la faire évoluer.
  \item Onglet \btn{Observations} : les mesures à saisir chaque jour en plus
    du poids (nombre ou texte, unité, obligatoire ou non), par exemple
    température ou consommation, puis \btn{Ajouter l'observation}.
  \item Onglet \btn{Points limites} : pour chaque critère, la condition
    (perte de poids en \%, score clinique total, score d'un item ou
    observation numérique), l'opérateur et le seuil, l'\btn{Action requise}
    (alerte seule, euthanasie, appel du vétérinaire, surveillance renforcée
    ou autre action) et le \btn{Message à l’expérimentateur}, puis
    \btn{Ajouter le point limite}.
\end{enumerate}
\end{procedure}
\capture[0.85\linewidth]{u_protocole}{Un plan / protocole : sévérité attendue et onglets grille, observations, points limites.}

\begin{procedure}{Créer une expérience et inclure des animaux}
\begin{enumerate}
  \item Dans le protocole, \btn{Expériences}, puis \btn{Nouvelle
    expérience} : numéro, titre et dates (comprises dans celles du
    protocole).
  \item \btn{Ajouter des animaux} : cherchez les animaux par numéro (choix
    multiple), ou \btn{Ajouter toute la cage}.
  \item Indiquez le \btn{Poids de référence (g)}, qui sert au calcul de la
    perte de poids, puis \btn{Ajouter à l'expérience}.
\end{enumerate}
\end{procedure}
\capture[0.85\linewidth]{u_experience}{Une expérience et ses animaux inclus.}
\begin{regle}
L'inclusion est refusée si l'autorisation de projet est expirée ou pas encore
valide, si l'espèce n'est pas autorisée, si l'animal n'a pas de lignée, ou si
le nombre d'animaux autorisés serait dépassé. Un animal compte une seule fois
par projet.
\end{regle}

\subsection{La page des expériences}
La page d'un protocole regroupe ses expériences.
\begin{itemize}
  \item \textbf{Gestion des points limites} (en tête) : les actions requises
    à traiter, avec une note de résolution ou l'enregistrement de
    l'euthanasie, puis les alertes sans action des 7 derniers jours et les
    actions déjà traitées.
  \item \btn{Animaux inclus} : état de chaque participation, bouton
    \btn{Suivi quotidien} et bouton \btn{Clôturer} pour mettre fin à la
    participation d'un animal qui reste dans l'établissement.
  \item \btn{Évaluations cliniques journalières} : tableau animaux × jours
    (score, poids ou variation de poids, code couleur, jours non saisis). Un
    clic sur une case ouvre la consultation du suivi, et sur un jour non
    saisi, la saisie.
  \item \textbf{Rapports} : rapport d'essai, rapport de plan / protocole
    (méthodologie, grille, critères de point limite) et \btn{Rapport (avec
    grilles)} de l'expérience (suivis, points limites, actes). Ce dernier
    s'exporte en CSV (données et actes). Tous sont imprimables.
\end{itemize}

\subsection{Le suivi journalier}
\begin{procedure}{Saisir le suivi d'un animal}
\begin{enumerate}
  \item Depuis l'expérience, ouvrez le suivi de l'animal. Pour une cage
    entière, utilisez le suivi groupé de la cage.
  \item Notez chaque item de la grille clinique (le total devient le score
    clinique), le poids et les observations.
  \item Confirmez les \textbf{actes expérimentaux} du jour : réalisé, ou non
    réalisé avec son motif.
  \item Enregistrez. Les points limites sont évalués aussitôt. Si un critère
    est atteint, une \textbf{action requise} est ouverte et les
    administrateurs sont prévenus.
\end{enumerate}
\end{procedure}
\capture[0.85\linewidth]{u_suivi}{Suivi journalier : action requise ouverte par un point limite, rappel des critères, saisie du jour.}

\begin{trace}
\begin{itemize}
  \item Un seul suivi par animal et par jour. Enregistrer à nouveau un jour
    déjà saisi est une \textbf{correction} : le \textbf{motif est
    obligatoire}, les valeurs remplacées sont conservées et affichées dans la
    consultation du suivi, et le jour est marqué « corrigé ».
  \item Une action d'euthanasie se solde en enregistrant la sortie de
    l'animal. Les autres actions se soldent par une note décrivant ce qui a
    été fait.
\end{itemize}
\end{trace}
\capture[0.85\linewidth]{u_suivi_correction}{Consultation d’un suivi corrigé : motif et valeurs remplacées.}

\subsection{Les actes expérimentaux}
Depuis une expérience, bouton \btn{Actes}, puis \btn{Planifier un acte}.

\begin{procedure}{Planifier un acte}
\begin{enumerate}
  \item \textbf{1. Acte} : \btn{Type d’acte} (injection, gavage, appât,
    anesthésie, analgésie, prélèvement, chirurgie, pesée de l'alimentation,
    autre), \btn{Nom de l’acte}, \btn{Voie d’administration},
    \btn{Substance / molécule}, \btn{Dose} et \btn{Unité de dose},
    \btn{Consignes}.
  \item Pour prélever le produit sur le stock, choisissez le
    \btn{Médicament du stock} et la quantité par animal. Avec une dose selon
    le poids (par exemple 75~mg/kg) et un médicament dont la concentration
    est connue, le volume est calculé pour chaque animal d'après son dernier
    poids.
  \item \textbf{2. Informations à demander lors de la réalisation}
    (facultatif) : par exemple \btn{+ Volume administré (mL)} ou
    \btn{+ Aliment restant (g)}, ou \btn{Ajouter une information} (nombre ou
    texte, unité, obligatoire ou non).
  \item \textbf{3. Dates} : \btn{Du}, \btn{Au} et \btn{tous les (jours)},
    puis \btn{Ajouter ces dates}. Recommencez pour d'autres périodes ;
    \btn{Effacer les dates} repart de zéro.
  \item \textbf{4. Animaux} : cochez tous les animaux, une cage entière ou
    les animaux un par un.
  \item \btn{Planifier l’acte}.
\end{enumerate}
\end{procedure}
\capture[0.85\linewidth]{u_planifier_acte}{Formulaire « Planifier un acte ».}

Chaque acte planifié s'affiche sous forme de grille animaux × dates :
$\checkmark$ réalisé, $\times$ non réalisé, $\circ$ à faire (sur fond jaune
s'il est en retard), --- annulé. Le survol d'une case donne l'heure, la
personne, le lot et les valeurs saisies. \btn{Annuler} retire les dates
restant à faire.
\capture[0.85\linewidth]{u_actes}{Un acte planifié : réalisé la veille, à faire aujourd’hui et demain.}

Le jour prévu, l'acte se confirme dans le suivi de l'animal ou de sa cage :
réalisé (avec les informations demandées, le lot et la quantité, le volume
selon le poids étant proposé), ou non réalisé avec son motif.
\begin{regle}
\begin{itemize}
  \item Chaque matin, les personnes concernées reçoivent les actes du jour et
    les actes en retard.
  \item Une personne dont la compétence n'est pas encore validée doit
    indiquer l'\textbf{encadrant présent} pour confirmer un acte réalisé.
  \item Une confirmation erronée se corrige avec un motif : ce qui avait été
    saisi est conservé, le médicament est recrédité et l'acte redevient à
    faire.
  \item Un médicament utilisé dans un acte expérimental est tracé comme une
    administration et débité du lot, sans ordonnance.
\end{itemize}
\end{regle}

\section{Pharmacie et traitements}\label{sec:pharmacie}
\menu{Animaux > Pharmacie, ordonnances et administrations}.

\subsection{Le stock de médicaments}
\begin{procedure}{Ajouter un médicament et recevoir un lot}
\begin{enumerate}
  \item \btn{Nouveau médicament} : \btn{Nom}, \btn{Présentation},
    \btn{Principe actif}, \btn{Unité de stock} (mL, comprimé…),
    \btn{Concentration} (par exemple 100~mg/mL, nécessaire au calcul d'une
    dose selon le poids), case \btn{Stupéfiant (registre spécial,
    inventaires)} et notes, puis \btn{Ajouter}.
  \item \btn{Réception d’un lot} : le médicament, \btn{N° lot},
    \btn{Péremption}, \btn{Quantité} reçue, date de \btn{Réception},
    \btn{Début d’utilisation} si le flacon est ouvert aussitôt,
    \btn{Emplacement} et \btn{Fournisseur}, puis \btn{Enregistrer le lot}.
\end{enumerate}
\end{procedure}
\capture[0.85\linewidth]{a_pharmacie}{Page Pharmacie : nouveau médicament et réception d’un lot.}

Le tableau \btn{Stock par lot} donne pour chaque lot sa péremption (en jaune
à moins de 30 jours, en rouge une fois expiré) et la quantité restante.
\begin{itemize}
  \item \btn{Ouvrir} : enregistrez la date d'ouverture du flacon (début
    d'utilisation). Une administration l'enregistre aussi automatiquement.
  \item \btn{Ajuster} : une quantité positive ou négative (\btn{+ / -}) et
    son \btn{Motif} (casse, inventaire, retour…), puis \btn{OK}.
\end{itemize}

\subsection{Les ordonnances}
\begin{procedure}{Enregistrer une ordonnance}
\begin{enumerate}
  \item \btn{Ordonnances} : sélectionnez les animaux (recherche à choix
    multiple, \btn{Ajouter la salle} pour tous les animaux d'une salle,
    \btn{Vider la sélection} pour recommencer).
  \item Le médicament et le nombre d'\btn{Administrations prévues par
    animal}.
  \item Le prescripteur, la référence de l'ordonnance, sa date, les dates de
    début et de fin du traitement.
  \item La dose et son unité, la voie, la fréquence (par exemple « 1 fois
    par jour ») et les \btn{Consignes}.
  \item \btn{Enregistrer l’ordonnance}. Elle apparaît dans
    l'\btn{Historique}, et les animaux traités sont signalés par une
    seringue sur la fiche de leur cage et dans la tournée.
\end{enumerate}
\end{procedure}
\capture[0.85\linewidth]{u_ordonnances}{Ordonnances : saisie et historique.}

\subsection{Les administrations}
\begin{procedure}{Administrer à toute une cage}
\begin{enumerate}
  \item Depuis la fiche de la cage ou la tournée, cliquez sur l'icône
    seringue : la page \btn{Administrations} s'ouvre pour cette cage.
  \item Pour chaque ordonnance en cours, les animaux encore à traiter sont
    affichés avec le nombre d'administrations déjà faites (par exemple
    « 2/3 »).
  \item Choisissez le \btn{Lot non périmé}, la \btn{Date / heure}, la
    \btn{Quantité / animal}, vérifiez la \btn{Voie} et ajoutez une note si
    besoin.
  \item \btn{Valider} : une administration est enregistrée pour chaque
    animal, et le lot est débité de la quantité totale.
\end{enumerate}
\end{procedure}
\begin{procedure}{Administrer à un seul animal}
\begin{enumerate}
  \item Section \btn{Administration individuelle} : le \btn{Traitement
    actif} (animal et médicament), puis le \btn{Lot non périmé}.
  \item \btn{Date / heure}, \btn{Quantité}, \btn{Voie} et \btn{Note}, puis
    \btn{Enregistrer l’administration}.
\end{enumerate}
\end{procedure}
\capture[0.85\linewidth]{u_administration_cage}{Administrations d’une cage : validation groupée et administration individuelle.}
\begin{regle}
\begin{itemize}
  \item Le lot doit être un lot du médicament prescrit ; un lot périmé est
    refusé, comme un lot dont le stock est insuffisant.
  \item Un animal ne reçoit pas plus d'administrations que le nombre prévu
    sur l'ordonnance.
  \item Les dernières administrations sont listées en bas de la page,
    celles des actes expérimentaux étant marquées \btn{Acte expérimental}.
\end{itemize}
\end{regle}
\begin{trace}
Une ordonnance enregistrée ne peut plus être modifiée ni supprimée. Une
administration enregistrée non plus : seule la correction d'un acte
expérimental annule l'administration correspondante, avec un motif, et
recrédite le lot. En cas d'erreur sur une administration d'ordonnance,
prévenez un administrateur ; l'écart de stock se corrige par un ajustement
motivé. Les mouvements des stupéfiants figurent dans le \btn{Registre des
stupéfiants}.
\end{trace}

\section{Terrariums et populations}\label{sec:terrariums}
\menu{Animaux > Terrariums et populations} : suivi des populations non
identifiées individuellement (rongeurs sauvages, amphibiens, insectes…). Le
tableau de bord donne pour chaque terrarium l'espèce ou la lignée, la salle,
l'estimation actuelle, le dernier recensement, la santé et le responsable.
La création, la modification et l'archivage d'un terrarium sont réservés aux
responsables BEA et aux administrateurs (voir le guide administrateur) ; toute personne habilitée enregistre le journal,
les traitements, les contrôles et les captures.

\capture[0.85\linewidth]{u_terrarium}{Page d’un terrarium : estimation, santé, délai d’attente et journal.}

La page d'un terrarium affiche trois encadrés : l'\textbf{estimation
actuelle} (dernier recensement corrigé des entrées et sorties depuis, et date
du prochain recensement, en jaune s'il est dû), la \textbf{santé} (dernier
contrôle, prochain contrôle, résultats en attente) et, le cas échéant, le
\textbf{délai d'attente} en cours. La courbe de population apparaît après
deux recensements.

\subsection{Recensements et journal}
\begin{procedure}{Ajouter une entrée au journal}
\begin{enumerate}
  \item Onglet \btn{Recensements et journal}, champ \btn{Entrée} :
    \begin{itemize}
      \item \btn{Recensement} : \btn{Estimation}, valeurs \btn{Bas} et
        \btn{Haut} de la fourchette et \btn{Méthode} (comptage visuel,
        comptage sur une zone échantillon, capture-recapture, estimation par
        volume ou poids) ;
      \item \btn{Introduction}, \btn{Mortalité constatée}, \btn{Capture} ou
        \btn{Transfert sortant} : le nombre d'animaux et le \btn{Stade}
        (œufs, larves, juvéniles, adultes, mixte) ;
      \item \btn{Note} : un texte libre (entretien, changement de litière…).
    \end{itemize}
  \item \btn{Date} (pas dans le futur), puis \btn{Ajouter au journal}.
\end{enumerate}
\end{procedure}
L'estimation affichée part du dernier recensement, ajoute les introductions
et retire les mortalités, captures et transferts enregistrés depuis.

\subsection{Traitements de la population}
\begin{procedure}{Enregistrer un traitement}
\begin{enumerate}
  \item Onglet \btn{Traitements} : \btn{Type} (antiparasitaire,
    antibiotique, antifongique, vitamines, désinfection du terrarium, autre)
    et \btn{Date / heure}.
  \item Soit un \btn{Médicament du stock (lot)} avec la \btn{Quantité}
    utilisée, débitée du lot ; soit une \btn{Substance (si hors stock)}.
  \item \btn{Dose}, \btn{Méthode} (bain, aliment, eau, pulvérisation…),
    \btn{Raison}, \btn{Prescripteur} et \btn{Ordonnance}.
  \item \btn{Délai d’attente jusqu’au} : la date jusqu'à laquelle les
    animaux ne doivent pas être utilisés en expérimentation.
  \item \btn{Enregistrer le traitement}.
\end{enumerate}
\end{procedure}
\capture[0.85\linewidth]{u_terrarium_traitements}{Onglet « Traitements » : un antiparasitaire avec délai d’attente.}

\subsection{Suivi sanitaire}
\begin{procedure}{Enregistrer un contrôle sanitaire}
\begin{enumerate}
  \item Onglet \btn{Suivi sanitaire}, champ \btn{Contrôle} :
    \begin{itemize}
      \item \btn{Observation sanitaire} : \btn{État général} normal ou
        anormal, \btn{Signes observés} (obligatoires si anormal) et
        \btn{Animaux malades vus} ;
      \item \btn{Paramètres d’ambiance} : \btn{Température (°C)},
        \btn{Humidité (\%)} et \btn{Autres paramètres} ;
      \item \btn{Dépistage sur prélèvements} : prélèvement, nombre
        d'animaux prélevés, agents recherchés, laboratoire et référence.
    \end{itemize}
  \item \btn{Mesures prises} (isolement, nettoyage, traitement, appel du
    vétérinaire…), puis \btn{Enregistrer le contrôle}.
  \item Pour un dépistage, le résultat est \btn{En attente}. À réception,
    choisissez \btn{Négatif} ou \btn{Positif}, la date et le détail, puis
    \btn{Enregistrer le résultat} (une seule fois).
\end{enumerate}
\end{procedure}
\capture[0.85\linewidth]{u_terrarium_sanitaire}{Onglet « Suivi sanitaire » : paramètres d’ambiance et dépistage en attente de résultat.}
\begin{regle}
Une observation anormale ou un dépistage positif est notifié aussitôt au
responsable du terrarium (à défaut, aux administrateurs). Au plus une fois
par semaine, le responsable est prévenu d'un recensement ou d'un contrôle
sanitaire en retard, selon les fréquences fixées pour le terrarium.
\end{regle}

\subsection{Captures}
Les animaux capturés deviennent des animaux identifiés : origine « capturé
dans une population de l'établissement » au registre des entrées, et ils
sortent de l'estimation de la population.
\begin{procedure}{Capturer des animaux}
\begin{enumerate}
  \item Onglet \btn{Captures} : \btn{Date}, \btn{Lignée},
    \btn{Date de naissance (estimée)}, \btn{Stade} et \btn{Salle des
    nouvelles cages}.
  \item \btn{Répartition par cage} : une ligne par cage de destination, avec
    le \btn{Nombre d’animaux} et le \btn{Sexe}. Laissez la cage vide pour
    créer une nouvelle cage d'hébergement, ou choisissez une cage existante.
    \btn{Remplissage rapide} crée d'un coup plusieurs cages d'un même
    nombre d'animaux ; \btn{Ajouter une cage} ajoute une ligne.
  \item Les numéros attribués s'affichent avant validation. Saisissez si
    possible le \textbf{poids individuel} de chacun.
  \item Pour une expérience : \btn{Inclure dans l’expérimentation} et, à
    défaut de poids individuels, un \btn{Poids de référence commun (g)}.
  \item \btn{Enregistrer la capture}. Les animaux capturés sont listés sous
    le formulaire avec leur cage, leur statut et leurs expériences.
\end{enumerate}
\end{procedure}
\capture[0.85\linewidth]{u_terrarium_captures}{Onglet « Captures » : répartition par cage et numéros attribués.}
\begin{regle}
Pendant un délai d'attente, aucune capture pour une expérimentation n'est
possible. Une capture hors expérimentation reste possible en cochant
\btn{Capturer pendant le délai d’attente (hors expérimentation)}.
\end{regle}

\section{Formation, tâches et rappels}
\begin{itemize}
  \item \menu{Paramètres > Mon livret de formation} : déclarez votre diplôme,
    vos formations réglementaires par fonction, vos compétences validées et
    votre formation continue, avec l'attestation PDF. Une entrée compte une
    fois validée par un administrateur ou un responsable BEA. Un bandeau vous
    prévient si votre formation continue arrive à échéance.
  \item \menu{Paramètres > Tâches et rappels} : tâches assignées, et rappels
    récurrents (quotidiens, hebdomadaires, mensuels) qui créent des tâches
    et envoient un e-mail.
\end{itemize}

\section{Les règles de traçabilité en bref}
\begin{trace}
\begin{itemize}
  \item Rien n'est effacé : registres sanitaire, des entrées et sorties, des
    stupéfiants, mouvements, journaux d'activité.
  \item Toute correction d'une donnée réglementaire demande un motif et
    conserve la valeur d'origine : suivis journaliers, confirmations d'actes,
    entrées et sorties du registre, administrations.
  \item Chaque saisie porte son auteur, sa date et son heure : travaillez
    avec votre propre compte, ne partagez jamais vos identifiants.
\end{itemize}
\end{trace}

\end{document}
